\chapter{Limitations and Conclusions}
\label{chap:conclusions}

\section{Limitations}
\label{sec:limitations}

The first limitation concerns what the validation shows. The cases of the suite
are vignettes that I wrote from the same reading of the criteria from which I
wrote the rules. Hence, a case that passes shows that the rules agree with that
reading, and the mutation test shows that the suite would notice most errors in
the rules. Neither shows that the reading is the right one, nor that the system
would reach the diagnosis of a clinician on real patients; this would require a
comparison on clinical records, which was not attempted. PESAD is therefore a
tool for exploring, checking and teaching the criteria of the manual; its result is not a diagnosis, as the bot and the web API state at
the start of every consultation.

Several limitations come from the way a consultation is conducted. The most
visible one is its length. A complete investigation poses several hundred
questions, since the knowledge base has 238 yes/no items, which
\code{unless\_absent/3} prunes only for the disorders whose core criteria have
already failed. This is acceptable in the terminal for a clinician, but rather tedious
in a chat, and it is the main reason why the \code{control} command and the
evaluation from test results (which can be restricted to some chapters or goals)
exist. A consultation by questions cannot yet be restricted to a chapter, and
the order of the questions is that of the rules, not that of an interview; a
screening question per disorder, posed before its criteria, would probably
shorten a consultation considerably. Furthermore, the duration of the symptoms is a single
answer, shared by all the disorders that state a duration (and, likewise, the
onset after a stressor and the time since a bereavement are single answers).
Hence, a patient with two disorders of different durations (say, panic attacks
for 4 months and generalized worry for 2 years) cannot be described, and the
answer given for one disorder is used for the other. Finally, the exact age is
known only in the evaluation from test results, whereas a consultation uses four
age bands. The thresholds that the manual places inside a band are thus applied
to the whole band; e.g., a 17-year-old is held to the 6 symptoms of ADHD that
apply below 17.

Exclusion criteria are, to a large extent, judgment questions. A criterion
such as ``the disturbance is not better explained by another mental disorder''
is a yes/no question for the clinician, and the system does not answer it from
its conclusions on the other goals. Hence, the hardest part of the differential
diagnosis is left to the user. Also, the list of differentials returned by the web
API is a reminder of what to rule out, rather than the result of having ruled it
out. The exceptions are the exclusions that the manual states in terms of other
criteria, which are encoded as such: e.g., social communication disorder is
separated from autism by the absence of restricted, repetitive behaviors, and
persistent tic disorder by never having met the criteria of Tourette's
disorder.

Being a variant of the certainty factors, the calculus shares their limits. The certainties of the rules ($0.95$ to $0.98$) were chosen by hand, and they
bound the certainty of a diagnosis from above. Similarly, the threshold of $0.5$ that
separates a present from an absent goal is a convention, not a calibrated value.
A certainty is a degree of belief, not a probability, and ``72\%'' for panic
disorder should not be read as a frequency. Each goal is the conclusion of a
single rule, so that the combination of the evidence of several rules for the
same conclusion, which MYCIN performed with a dedicated
formula~\cite{shortliffe-buchanan}, is not needed. On the other hand, for the
polythetic disorders the choice of the t-norm has an effect that is larger than
one would wish (table~\ref{tab:uncertainty-cases}), and which only the minimum
avoids. As for coverage, the \emph{other specified} and \emph{unspecified}
categories, a few neurodevelopmental disorders (global developmental delay,
speech sound disorder, childhood-onset fluency disorder, stereotypic movement
disorder), the panic attack specifier and the specifiers of autism on
associated conditions and catatonia are not modeled. The ICD-11 codes are
recorded as metadata, without the ICD-11 criteria, which differ from those of
the DSM-5-TR in several places.


\section{Conclusions}

I described PESAD, a rule-based expert system that encodes the DSM-5-TR
criteria of 92 codable goals of the neurodevelopmental, anxiety,
obsessive-compulsive and trauma-related chapters, with their exclusions,
specifiers and ICD-10-CM and ICD-11 codes. It reasons under uncertainty with
certainty factors, i.e., with the minimum t-norm, or with one of five
alternative methods. The same engine serves a terminal shell, a Telegram bot
and a stateless web API, which also evaluates the knowledge base from test
results with lower and upper certainty bounds and the list of the missing data.
A validation suite of 425 cases per method, covering each goal with a typical,
a subthreshold, an exclusion and a differential case, passes under all six
methods, and it detects 128 of 145 errors injected into the knowledge base. Comparing the methods supports the minimum as the default, since the other
t-norms make the certainty of a diagnosis depend on the number of its criteria.
Open questions are mostly practical: whether the consultation can be
shortened by screening questions without losing comorbidities, and whether the
exclusions can be answered (at least in part) from the conclusions of the
system. Above all, it remains to be seen whether the conclusions of the system
agree with those of clinicians on real cases, a question that a knowledge base
built from a manual cannot settle.
